% GENERATED FILE. Edit canonical lessons, then run npm run generate:books.
% canonical-source-hash: fnv1a64:d4c64a5e309e20e0
% canonical-lessons: FA-C13-aseman, FA-C13-khorshid, FA-C13-mah, FA-C13-setare, FA-C13-baran

\chapter{Sky, Sun, Moon, Star, Rain}
\label{ch:fa-sky-words}
\hlchaptermodality{13}

\begin{chapteropening}
\textbf{By the end of this chapter:} \emph{I can name the sky, sun, moon, star, and rain in Persian, and I can say which of those words has no established English cousin.}

The chapter closes on bârân, a word with no English cousin, then runs all five sky words back to back and rebuilds the careful wellbeing question the weather now follows in ordinary small talk.
\end{chapteropening}

\section[âsemân]{\fa{آسمان} — sky}

\label{lesson:FA-C13-aseman}



\begin{quote}
Say \textbf{\fa{کتاب}} (\emph{ketâb}), and which language it came from. This chapter turns from things in the house to the world above it, starting with the biggest word of all.
\end{quote}

\subsection*{You'll want to know first — one word}
\begin{quote}
\textbf{\fa{آسمان}} — \emph{âsemân} — \textbf{sky}
\end{quote}

Five letters, all already yours: \textbf{\fa{آ} \fa{س} \fa{م} \fa{ا} \fa{ن}} — alef with \emph{madde}, \emph{s}, \emph{m}, long \emph{â}, \emph{n}.

\begin{cousinweb}
\textbf{\fa{آسمان}} continues Old Persian \textbf{asmānam}, from Indo-European *\emph{h\textsubscript{2}é\'{k}mō}, \textquotedblleft{}stone.\textquotedblright{} Sanskrit kept the plain sense as \textbf{aśman}, \textquotedblleft{}stone,\textquotedblright{} while Avestan's cousin word already carried both meanings, \textquotedblleft{}stone\textquotedblright{} and \textquotedblleft{}sky,\textquotedblright{} at once. Picturing the sky as a stone dome overhead was a widespread Indo-European habit, not a Persian invention — one old word doing two jobs, the way \textbf{\fa{زبان}} later did for tongue and language.
\end{cousinweb}

\subsection*{Your turn}
\begin{itemize}
\raggedright
  \item \emph{Say it:} \textbf{âsemân} — sky
  \item \emph{Connect:} \textbf{âsemân} $\leftarrow$ *\emph{h\textsubscript{2}é\'{k}mō} $\to$ Sanskrit \textbf{aśman}, \textquotedblleft{}stone\textquotedblright{}
  \item \emph{Name:} the old Indo-European habit \textbf{\fa{آسمان}} preserves — picturing the sky as stone
  \item \emph{Retrieve:} \textbf{\fa{زبان}} — tongue and language, with its letter \textbf{\fa{ز}}
\end{itemize}

\begin{tcolorbox}[breakable,colback=teal!4,colframe=teal!35!black,title={Before you move on}]
What does \textbf{\fa{آسمان}} mean? (\textbf{Sky.}) What did its root originally mean, and which Sanskrit word still keeps that older sense? (\textbf{Stone}; \textbf{aśman}.)

Source: \href{https://en.wiktionary.org/wiki/\%D8\%A2\%D8\%B3\%D9\%85\%D8\%A7\%D9\%86}{Wiktionary: \fa{آسمان}}.
\end{tcolorbox}
\section[khorshid]{\fa{خورشید} — sun, a word built from two old pieces}

\label{lesson:FA-C13-khorshid}



\begin{quote}
Say \textbf{\fa{آسمان}} (\emph{âsemân}), and the root sense it once carried. Look up, and this lesson names what is in it.
\end{quote}

\subsection*{You'll want to know first — one word}
\begin{quote}
\textbf{\fa{خورشید}} — \emph{khorshid} — \textbf{sun}
\end{quote}

Six letters, all already yours: \textbf{\fa{خ} \fa{و} \fa{ر} \fa{ش} \fa{ی} \fa{د}} — \emph{kh}, \textbf{\fa{و}} carrying its learned \emph{o} value, \emph{r}, \emph{sh}, long \emph{i}, \emph{d}.

\begin{cousinweb}
\textbf{\fa{خورشید}} is a compound, not a single old root. Its first half, \textbf{\fa{خور}} \emph{khor}, \textquotedblleft{}sun,\textquotedblright{} descends from the ancient Indo-European word for the sun, *\emph{sóh\textsubscript{2}wl} — the same distant root behind English \textbf{sun}, Latin \textbf{sol} ($\to$ \textbf{solar}), Greek \textbf{h\'{\={e}}lios} ($\to$ \textbf{heliotrope}), and Sanskrit \textbf{sūrya}, though it reached Persian by a completely different branch of the family. Its second half, \textbf{\fa{شید}} \emph{shid}, \textquotedblleft{}shining,\textquotedblright{} is an unrelated word that Avestan had already paired with \textbf{\fa{خور}} long before Persian existed.
\end{cousinweb}

\subsection*{Your turn}
\begin{itemize}
\raggedright
  \item \emph{Say it:} \textbf{khorshid} — sun
  \item \emph{Split:} \textbf{\fa{خور}}, sun — \textbf{\fa{شید}}, shining
  \item \emph{Connect:} \textbf{\fa{خور}} $\leftarrow$ *\emph{sóh\textsubscript{2}wl} $\to$ English \textbf{sun}, Latin \textbf{sol}, Greek \textbf{h\'{\={e}}lios}, Sanskrit \textbf{sūrya}
  \item \emph{Run through:} \textbf{âsemân, khorshid} — sky, sun
\end{itemize}

\begin{tcolorbox}[breakable,colback=teal!4,colframe=teal!35!black,title={Before you move on}]
What does \textbf{\fa{خورشید}} mean, and what are its two parts? (\textbf{Sun}; \textbf{\fa{خور}} \textquotedblleft{}sun\textquotedblright{} plus \textbf{\fa{شید}} \textquotedblleft{}shining.\textquotedblright{}) Which English word shares \textbf{\fa{خور}}'s distant root? (\textbf{Sun.})

Source: \href{https://en.wiktionary.org/wiki/\%D8\%AE\%D9\%88\%D8\%B1\%D8\%B4\%DB\%8C\%D8\%AF}{Wiktionary: \fa{خورشید}}.
\end{tcolorbox}
\section[mâh]{\fa{ماه} — moon, and month}

\label{lesson:FA-C13-mah}



\begin{quote}
Say \textbf{\fa{خورشید}} (\emph{khorshid}), and the two old pieces fused inside it. The sky's other light gives Persian one more word doing double duty.
\end{quote}

\subsection*{You'll want to know first — one word}
\begin{quote}
\textbf{\fa{ماه}} — \emph{mâh} — \textbf{moon}, and also \textbf{month}
\end{quote}

Three letters, all already yours: \textbf{\fa{م} \fa{ا} \fa{ه}} — \emph{m}, long \emph{â}, \emph{h}.

\begin{cousinweb}
\textbf{\fa{ماه}} continues Old Persian \textbf{māha-}, from Indo-European *\emph{m\'{\={e}}h\textsubscript{1}ns}, \textquotedblleft{}moon.\textquotedblright{} The same root split into two related English words, \textbf{moon} and \textbf{month} — a month was, for most of human history, one full cycle of the moon — and Persian's single word still holds both senses together, the way \textbf{\fa{زبان}} later holds \textquotedblleft{}tongue\textquotedblright{} and \textquotedblleft{}language\textquotedblright{} in one word. Latin kept the root as \textbf{mensis} ($\to$ \textbf{month}, \textbf{menstrual}), Greek as \textbf{mēn}, and Sanskrit as \textbf{māsa}.
\end{cousinweb}

\subsection*{Your turn}
\begin{itemize}
\raggedright
  \item \emph{Say it:} \textbf{mâh} — moon, and month
  \item \emph{Connect:} \textbf{mâh} $\leftarrow$ *\emph{m\'{\={e}}h\textsubscript{1}ns} $\to$ English \textbf{moon} and \textbf{month}, Latin \textbf{mensis}, Sanskrit \textbf{māsa}
  \item \emph{Run through:} \textbf{âsemân, khorshid, mâh} — sky, sun, moon
\end{itemize}

\begin{tcolorbox}[breakable,colback=teal!4,colframe=teal!35!black,title={Before you move on}]
What two English words does \textbf{\fa{ماه}}'s root split into? (\textbf{Moon} and \textbf{month}.) Which earlier Persian word did the same double duty for tongue and language? (\textbf{\fa{زبان}}.)

Source: \href{https://en.wiktionary.org/wiki/\%D9\%85\%D8\%A7\%D9\%87}{Wiktionary: \fa{ماه}}.
\end{tcolorbox}
\section[setâre]{\fa{ستاره} — star}

\label{lesson:FA-C13-setare}



\begin{quote}
Say \textbf{\fa{ماه}}, and the two English words its root split into. One more light in the sky, and a new small spelling habit.
\end{quote}

\subsection*{You'll want to know first — one word, one small spelling habit}
\begin{quote}
\textbf{\fa{ستاره}} — \emph{setâre} — \textbf{star}
\end{quote}

Five letters, all already yours: \textbf{\fa{س} \fa{ت} \fa{ا} \fa{ر} \fa{ه}} — \emph{s}, \emph{t}, long \emph{â}, \emph{r}, and a final \textbf{\fa{ه}} read as the vowel \emph{-e} rather than as a consonant — the first time this track has met \textbf{\fa{ه}} doing that particular job.

\begin{cousinweb}
\textbf{\fa{ستاره}} continues Old Persian \textbf{star-}, from Indo-European *\emph{h\textsubscript{2}st\'{\={e}}r}, \textquotedblleft{}star.\textquotedblright{} Few Indo-European roots survive this cleanly across the whole family: English \textbf{star}, Latin \textbf{stella} ($\to$ \textbf{stellar}, \textbf{constellation}), Greek \textbf{astron} ($\to$ \textbf{astronomy}, \textbf{disaster} — \textquotedblleft{}ill-starred\textquotedblright{}), and Persian's own \textbf{\fa{ستاره}} all continue the identical ancient word, alongside \textbf{\fa{در}}'s door and \textbf{\fa{ماه}}'s moon as the tranche's clearest cousins.
\end{cousinweb}

\subsection*{Your turn}
\begin{itemize}
\raggedright
  \item \emph{Say it:} \textbf{setâre} — star
  \item \emph{Connect:} \textbf{setâre} $\leftarrow$ *\emph{h\textsubscript{2}st\'{\={e}}r} $\to$ English \textbf{star}, Latin \textbf{stella}, Greek \textbf{astron}
  \item \emph{Read it:} the final \textbf{\fa{ه}} in \textbf{\fa{ستاره}} as \emph{-e}, not as a consonant
  \item \emph{Run through:} \textbf{âsemân, khorshid, mâh, setâre} — sky, sun, moon, star
  \item \emph{Retrieve:} \textbf{\fa{در}} — door, and its own well-attested cousin set
\end{itemize}

\begin{tcolorbox}[breakable,colback=teal!4,colframe=teal!35!black,title={Before you move on}]
What does \textbf{\fa{ستاره}} mean? (\textbf{Star.}) How is its final \textbf{\fa{ه}} read? (As the vowel \textbf{-e}, not a consonant.) Name two other languages whose word for star shares the same root. (\textbf{English} \emph{star}, \textbf{Latin} \emph{stella}, or \textbf{Greek} \emph{astron}.)

Source: \href{https://en.wiktionary.org/wiki/\%D8\%B3\%D8\%AA\%D8\%A7\%D8\%B1\%D9\%87}{Wiktionary: \fa{ستاره}}.
\end{tcolorbox}
\section[bârân]{\fa{باران} — rain, and the wellbeing question this chapter answers differently}

\label{lesson:FA-C13-baran}



\begin{quote}
Say \textbf{\fa{ستاره}} (\emph{setâre}), and its root's three attested cousins. Not every sky word gets to keep company like that.
\end{quote}

\subsection*{You'll want to know first — one word}
\begin{quote}
\textbf{\fa{باران}} — \emph{bârân} — \textbf{rain}
\end{quote}

Five letters, all already yours: \textbf{\fa{ب} \fa{ا} \fa{ر} \fa{ا} \fa{ن}} — \emph{b}, long \emph{â}, \emph{r}, long \emph{â}, \emph{n}.

\begin{cousinweb}
\textbf{\fa{باران}} continues Middle Persian \textbf{wārān}, from Indo-Iranian *\emph{wáHr}, \textquotedblleft{}water, rain,\textquotedblright{} itself from Indo-European *\emph{weh\textsubscript{1}r-}, \textquotedblleft{}water.\textquotedblright{} Its cousins stayed close to home: Avestan \textbf{vāra}, \textquotedblleft{}rain,\textquotedblright{} and Sanskrit \textbf{vār}, \textquotedblleft{}water,\textquotedblright{} share the identical root. English \textbf{rain} is not part of this family at all — its own origin is disputed and confined to the Germanic languages, with no secure Indo-European cousin beyond them — so, as with \textbf{\fa{رفتن}} and \textbf{\fa{گفتن}} far earlier, the honest answer here is that no English cousin is claimed.
\end{cousinweb}

\begin{grammarlens}[title={Grammar Lens: what small talk about the sky follows}]
Weather small talk in Persian commonly follows the wellbeing exchange, the same slot English fills right after \textquotedblleft{}how are you?\textquotedblright{} Rebuild the question: \textbf{\fa{حال} \fa{شما} \fa{چطور} \fa{است؟}}, careful \textbf{ast}, not a contraction, then its reply, \textbf{\fa{خوبم،} \fa{ممنون}} — a shape a sky word can now attach to.
\end{grammarlens}

\subsection*{Your turn — every sky word, then the exchange it follows}
\begin{itemize}
\raggedright
  \item \emph{Say it:} \textbf{bârân} — rain
  \item \emph{Connect:} \textbf{bârân} $\leftarrow$ *\emph{weh\textsubscript{1}r-} $\to$ Sanskrit \textbf{vār}, Avestan \textbf{vāra} — no English cousin
  \item \emph{Run through:} \textbf{âsemân, khorshid, mâh, setâre, bârân} — all five sky words
  \item \emph{Rebuild:} \textbf{\fa{حال} \fa{شما} \fa{چطور} \fa{است؟}} and its reply \textbf{\fa{خوبم،} \fa{ممنون}}
\end{itemize}

\begin{tcolorbox}[breakable,colback=teal!4,colframe=teal!35!black,title={Before you move on}]
What does \textbf{\fa{باران}} mean? (\textbf{Rain.}) Does it share a root with English \textbf{rain}? (\textbf{No.}) Which two Indo-Iranian languages keep the closest cousins? (\textbf{Avestan}, \textbf{Sanskrit}.) Name all five sky words. (\textbf{Âsemân, khorshid, mâh, setâre, bârân}.)

Source: \href{https://en.wiktionary.org/wiki/\%D8\%A8\%D8\%A7\%D8\%B1\%D8\%A7\%D9\%86}{Wiktionary: \fa{باران}}.
\end{tcolorbox}
